\section{Scope, arithmetic target, and principal results}\label{sec:scope}

The problem is to construct one continuous linear source law
\begin{equation}\label{eq:goal}
J:\cD^0\longrightarrow\cH_{\rm YM},\qquad
\ip{Jf}{Jg}_{\rm OS}=Q(f,g),
\end{equation}
where the physical model, state, observable prescription and completion
are fixed for all compact arithmetic supports. The sources must arise
from independently specified YM observables or relations. An arbitrary
Hilbert embedding of a form already assumed positive does not solve this
problem. Assuming OS positivity, \eqref{eq:goal} implies RH by Weil's
criterion. A nonconstructive source-occurrence theorem is an acceptable
route, provided its feasibility hypotheses are independently established.

Most results below concern pure SU(2) Wilson theory on a spatial
$6^3$ torus and physical-time slices $-2,-1,0,1,2$, with open outer
boundaries, coupling $\beta=8/5$, and theta angle zero. The boundary
state is a finite-slab state, not an identified continuum vacuum.
The two-dimensional and thermal calculations in the appendices are
explicit comparisons, not changes to this main model.

Full OS reconstruction, continuum YM existence, and a physical mass gap
\cite{jaffe-witten}
may be assumed where explicitly stated. Local phase-space compactness
is a further optional physical assumption. These hypotheses have useful
roles in positivity and source compactness, but no argument here derives
the missing arithmetic source identity from them. The finite-state
weighted electric Laplacian is never identified with the physical
continuum Hamiltonian. Arithmetic translation, physical time, contour
capacity, and the representation scale are distinct parameters.

\subsection{Conventions and the complete form}
Let $\cD=C_c^\infty(\R)$, with its usual compact-support test topology,
and use inner products antilinear in the first entry. Put
\begin{equation}
\widehat f(\tau)=\int_\R e^{-\ii\tau x}f(x)\dd x,
\qquad U_tf(x)=f(x-t),\qquad
\T=-\partial_x^2+\tfrac14.
\end{equation}
The pole-neutral space is
\begin{equation}\label{eq:pole-neutral}
\cD^0=\left\{f\in\cD:
\int e^{x/2}f(x)\dd x=\int e^{-x/2}f(x)\dd x=0\right\}.
\end{equation}
For $r>0$, set
\begin{equation}
n_\Gamma(r)=\frac{e^{-r/2}}{1-e^{-2r}},\qquad
m_+(\tau)=\Ree\psi(\tfrac14+\ii\tau/2)-\log\pi,
\qquad \psi=\Gamma'/\Gamma.
\end{equation}
The full quadratic form is
\begin{align}\label{eq:full-weil}
Q[f]&=Q_\infty[f]+P[f]
-2\sum_{a\ge2}\frac{\Lambda(a)}{\sqrt a}
          \Ree\ip{f}{U_{\log a}f}_2,\\
Q_\infty[f]&=\frac1{2\pi}\int_\R m_+(\tau)|\widehat f(\tau)|^2\dd\tau,
\label{eq:arch-fourier}\\
P[f]&=2\left|\int f(x)\cosh(x/2)\dd x\right|^2
      -2\left|\int f(x)\sinh(x/2)\dd x\right|^2.
\end{align}
Here $\Lambda(p^k)=\log p$ and is zero at other integers.
Polarization defines $Q(f,g)$. The prime sum is finite for each compact
test pair, because sufficiently large shifts have disjoint supports.
The pole term vanishes on $\cD^0$. On that space $Q$ is stationary
and invariant under reflection $f(x)\mapsto f(-x)$.

Equivalently,
\begin{equation}\label{eq:arch-difference}
Q_\infty[f]=m_+(0)\norm f_2^2+
\frac12\iint_{\R^2}|f(x)-f(y)|^2n_\Gamma(|x-y|)\dd x\dd y.
\end{equation}
The difference makes the diagonal singularity integrable. Its completed
contact and asymptotic are
\begin{equation}\label{eq:contact-growth}
m_+(0)=-\gamma_{\!E}-\frac\pi2-3\log2-\log\pi,
\qquad
m_+(\tau)=\log\frac{|\tau|}{2\pi}+O(|\tau|^{-2}).
\end{equation}
These formulas follow from the digamma integral and asymptotics
\cite{dlmf}. In particular, the off-diagonal kernel alone does not fix
the required form. The normalization of \eqref{eq:full-weil} and the
sufficiency of positivity on \eqref{eq:pole-neutral} are the restricted
Weil criterion in \cite[Appendices B--C]{cc}.

\begin{lemma}[Exact removal of the pole moments]\label{lem:T}
The map $\T:\cD\to\cD^0$ is a topological isomorphism. Its inverse is
\begin{equation}
(\T^{-1}f)(x)=\int_\R e^{-|x-y|/2}f(y)\dd y.
\end{equation}
The inverse preserves an enclosing support interval, although it need
not preserve holes in a disconnected support.
\end{lemma}
\begin{proof}
The displayed kernel is the fundamental solution of
$-\partial^2+1/4$. Its two exterior tails vanish precisely by the two
moments. The resulting function is smooth, since its equation is
$h''=h/4-f$, and is supported in the convex hull of $\supp f$.
Integration by parts proves $\T\cD\subset\cD^0$; compactly supported
solutions of $\T h=0$ vanish. The integral formula and differential
equation give continuity on each fixed-support test space, hence the
claimed test-space isomorphism.
\end{proof}

\begin{proposition}[A necessary ultraviolet test]\label{prop:highfreq}
For nonzero $h\in\cD$, the exactly pole-neutral functions
\begin{equation}\label{eq:modulation}
f_\lambda=\lambda^{-2}\T(e^{\ii\lambda x}h)
=e^{\ii\lambda x}\left(h-\frac{2\ii h'}\lambda+
                     \frac{h/4-h''}{\lambda^2}\right)
\end{equation}
have uniformly bounded ordinary $L^1,L^2,L^\infty$ norms and satisfy
\begin{equation}\label{eq:highfreq}
Q[f_\lambda]=\norm h_2^2\log|\lambda|+O_h(1)
\quad (|\lambda|\to\infty).
\end{equation}
Consequently no source law locally bounded solely in one of those
ordinary input norms realizes $Q$. In particular this excludes
$Jf=\int f(x)A(x)\dd x$ when $A$ is a Hilbert-valued function with
locally integrable norm.
\end{proposition}
\begin{proof}
The unmodulated profile in \eqref{eq:modulation} converges smoothly to h
on a fixed support and has uniform Schwartz Fourier estimates.
In \eqref{eq:arch-fourier} split at $|\tau-\lambda|=|\lambda|/2$.
Equation \eqref{eq:contact-growth} gives the leading logarithm on the
first region and the Schwartz tail bounds the second. The finitely many
prime terms on the fixed support remain bounded by Cauchy--Schwarz.
For the stated integral law,
$\norm{Jf_\lambda}\le\norm{f_\lambda}_\infty
\int_{\supp h}\norm{A(x)}\dd x$, contradicting \eqref{eq:highfreq}.
\end{proof}

\subsection{What has been established}
The affirmative identities in this manuscript are the exact reflected
boundary norm, two actual generalized-source constructions with
nonarithmetic pairings, the character prime mixed limit, the complete
signed archimedean/prime limit, and several sufficient occurrence and
determination theorems. The exclusions identify where these identities
fail to become a positive Weil-source law. The last criterion reduces
the target to one stationary physical correlation, but that correlation
is not obtained from a YM observable here. Classical conductor-operator
and explicit-formula mathematics is credited to its sources; no claim
of novelty is made for those ingredients or for the present deductions
without further literature and independent mathematical review.
